% =====================================================================
%  1bat-ud2-6.tex  —  HORA 6: Sumar i restar pressupostos
%  (sense preàmbul: s'inclou des de main.tex)
% =====================================================================
\horatitol{Sessió 6 — Sumar i restar pressupostos}%
{Identificar termes semblants i operar amb polinomis: suma i resta, controlant
el signe en cada pas.}

\subsection{Idea clau}

A la sessió anterior vam calcular valors numèrics. Avui combinem dos pressupostos
expressats com a polinomis per obtenir-ne un de sol. La clau és que
\textbf{només es poden ajuntar els termes del mateix tipus}: els termes en $x^2$
amb els de $x^2$, els de $x$ amb els de $x$, i els nombres sols amb els nombres
sols. Quan restem, el parany és oblidar de canviar \emph{tots} els signes del
segon polinomi.

\begin{idea}
\textbf{Termes semblants:} dos termes són semblants si tenen la mateixa part
literal (la mateixa lletra elevada al mateix grau).
\begin{itemize}[leftmargin=1.4em, itemsep=2pt]
  \item $3x^2$ i $-5x^2$ són semblants; $3x^2$ i $3x$ \emph{no} ho són.
\end{itemize}
\textbf{Suma de polinomis:} s'agrupen i s'operen els termes semblants:
\[
(ax^2 + bx + c) + (dx^2 + ex + f) = (a+d)x^2 + (b+e)x + (c+f).
\]
\textbf{Resta de polinomis:} es canvien \emph{tots} els signes del segon polinomi
i després es suma:
\[
(ax^2+bx+c)-(dx^2+ex+f) = ax^2+bx+c \mathbf{-}dx^2 \mathbf{-}ex \mathbf{-}f.
\]
\textbf{Compte amb el signe!} En la resta, el signe «$-$» davant del parèntesi
afecta tots els termes de dins, no només el primer.
\end{idea}

\subsection{Exemples}

\begin{exemple}[combinar dos pressupostos parcials d'un servei]
Un projecte d'atenció domiciliària té dos pressupostos en funció del nombre de
persones ateses~$x$:
\begin{itemize}[leftmargin=1.4em, itemsep=2pt]
  \item Cost de personal: $P(x) = 20x + 800$ (en euros)
  \item Cost de materials: $M(x) = 5x + 150$ (en euros)
\end{itemize}
El pressupost total és:
\begin{align*}
T(x) &= P(x) + M(x) \\
     &= (20x + 800) + (5x + 150) \\
     &= (20+5)x + (800+150) \\
     &= 25x + 950.
\end{align*}
Interpretació: cada persona atesa costa $25\eur$ (sumant personal i materials)
i hi ha un cost fix de $950\eur$.
\end{exemple}

\begin{exemple}[restar per trobar el cost addicional d'un nou servei]
Una entitat ja gestiona un servei social amb cost $A(x) = 3x^2 + 50x + 400$ i
n'obre un de nou amb cost $B(x) = x^2 + 30x + 200$. Quant costa \emph{més}
el primer que el segon?
\begin{align*}
A(x) - B(x) &= (3x^2 + 50x + 400) - (x^2 + 30x + 200) \\
             &= 3x^2 + 50x + 400 - x^2 - 30x - 200 \\
             &= (3-1)x^2 + (50-30)x + (400-200) \\
             &= 2x^2 + 20x + 200.
\end{align*}
Observació: en treure el parèntesi negatiu, \textbf{tots} els signes canvien:
$-x^2$, $-30x$, $-200$.
\end{exemple}

\subsection{Exercicis resolts}

\begin{exresolt}[pressupost total d'un programa social]
Un programa de serveis comunitaris té un cost de coordinació
$C(x) = 2x^2 + 100x + 500$ i un cost d'execució $E(x) = x^2 + 60x + 300$,
on $x$ és el nombre de participants. Calcula el pressupost total $T(x) = C(x) + E(x)$.
\solucio
\begin{align*}
T(x) &= (2x^2 + 100x + 500) + (x^2 + 60x + 300) \\
     &= 2x^2 + x^2 + 100x + 60x + 500 + 300
       &&\text{(agrupo per tipus)} \\
     &= 3x^2 + 160x + 800.
       &&\text{(sumo semblants)}
\end{align*}
\resposta{$T(x) = 3x^2 + 160x + 800\eur$}
\end{exresolt}

\begin{exresolt}[reducció de pressupost per retallada]
Un casal de barri tenia un pressupost $P(x) = 4x^2 + 80x + \num{1200}$ i,
a causa d'una retallada, el nou pressupost és $Q(x) = x^2 + 30x + 400$.
Calcula la reducció $R(x) = P(x) - Q(x)$.
\solucio
\begin{align*}
R(x) &= (4x^2 + 80x + \num{1200}) - (x^2 + 30x + 400) \\
     &= 4x^2 + 80x + \num{1200} - x^2 - 30x - 400
       &&\text{(canvio tots els signes del segon)} \\
     &= (4-1)x^2 + (80-30)x + (\num{1200}-400) \\
     &= 3x^2 + 50x + 800.
\end{align*}
\resposta{La reducció és $R(x) = 3x^2 + 50x + 800\eur$.}
\end{exresolt}

\subsection{Exercicis per fer}

\begin{exercicis}
\begin{exlist}
  \item Calcula les sumes de polinomis següents:
    \begin{enumerate}[label=\alph*), leftmargin=1.6em, topsep=2pt]
      \item $(3x + 7) + (5x + 2)$
      \item $(x^2 + 4x - 1) + (2x^2 - x + 6)$
    \end{enumerate}

  \item Calcula les restes de polinomis següents:
    \begin{enumerate}[label=\alph*), leftmargin=1.6em, topsep=2pt]
      \item $(8x + 15) - (3x + 9)$
      \item $(4x^2 + 7x + 10) - (x^2 + 7x + 3)$
    \end{enumerate}

  \item Un ajuntament finança dos serveis d'atenció a persones grans:
    \begin{itemize}[leftmargin=1.4em, itemsep=2pt, topsep=2pt]
      \item Teleassistència: $T(x) = 8x + 600$
      \item Servei d'àpats: $S(x) = 12x + 400$
    \end{itemize}
    on $x$ és el nombre d'usuaris. Calcula el cost total $T(x)+S(x)$ i
    interpreta el coeficient de $x$ i el terme independent del resultat.

  \item Una entitat rep una subvenció de $G(x) = 2x^2 + 50x + \num{1000}$
        i té uns costos de $D(x) = x^2 + 30x + 600$, on $x$ és el nombre
        d'activitats. Calcula el romanent $G(x) - D(x)$ (el que sobra un cop
        pagats els costos).

  \item \textbf{(Compte amb el signe)} Calcula $(5x^2 - 3x + 8) - (5x^2 + 2x - 1)$.
        Abans de resoldre-ho, encercla el signe de cadascun dels tres termes
        del segon polinomi que \emph{canviaran} de signe en treure el parèntesi.

  \item \textbf{(Raonament)} Dues entitats gestionen conjuntament un projecte.
        La primera té el pressupost $A(x) = 3x^2 + 40x + 700$ i la segona
        $B(x) = 3x^2 + 40x + 700$ (exactament el mateix). Sense fer cap càlcul,
        digues quin és el resultat de $A(x) - B(x)$ i explica, en el context
        dels pressupostos, \emph{per què} té sentit aquest resultat.
\end{exlist}
\end{exercicis}

% ---------------------------------------------------------------------
\fullsolucions{Sessió 6}
\begin{solucions}
\begin{sollist}
  \item (a) $8x+9$ \quad (b) $3x^2+3x+5$
  \item (a) $5x+6$ \quad (b) $3x^2+7$
  \item $T(x)+S(x) = 20x+\num{1000}$. El coeficient $20$ és el cost per usuari
        addicional (suma de $8+12$); el terme $\num{1000}$ és el cost fix total
        dels dos serveis.
  \item $G(x)-D(x) = x^2+20x+400$
  \item $(5x^2-3x+8)-(5x^2+2x-1) = 5x^2-3x+8-5x^2-2x+1 = -5x+9$
  \item $A(x)-B(x)=0$. Dos pressupostos idèntics es cancel·len completament:
        la diferència entre les despeses de les dues entitats és zero, cosa que
        indica que contribueixen exactament el mateix al projecte.
\end{sollist}
\end{solucions}
